\chapter{MOTOR CONTROL AND TRACKING LOOPS}
\label{ch:motor}

This chapter follows a velocity command from \topic{/cmd\_vel} down to the
voltage across each motor, models the motors, lists every feedback and
tracking loop on the robot with its rate, and then proposes the two loops the
robot is missing: a closed-loop wheel-speed controller and a yaw-rate
controller. Sections marked \emph{proposed} describe designs that are not yet
in the code.

\section{THE ACTUATOR CHAIN}
\label{sec:actuator}

\begin{figure}[htbp]
\centering
\begin{tikzpicture}[
  box/.style={draw,rounded corners=2pt,align=center,font=\small,
              minimum height=10mm,inner sep=3pt,text width=22mm},
  arr/.style={-{Stealth[length=2mm]},thick}]
\node[box] (cmd) {\topic{/cmd\_vel}\\$(v,\omega)$};
\node[box,right=6mm of cmd] (side) {Side speeds\\\eqref{eq:sidespeed}};
\node[box,right=6mm of side] (duty) {Duty map\\\eqref{eq:duty}};
\node[box,right=6mm of duty,text width=27mm] (rc) {{\footnotesize\texttt{rc\_motor\_set}}\\dir GPIOs + PWM};
\node[box,below=8mm of rc,text width=27mm] (hb) {TB6612FNG\\H-bridge};
\node[box,left=12mm of hb] (mot) {DC gear\\motor ($\times2$)};
\node[box,left=6mm of mot] (wheel) {Wheel and\\floor};
\node[box,left=6mm of wheel] (enc) {Hall encoder\\(odometry only)};
\draw[arr] (cmd) -- (side);
\draw[arr] (side) -- (duty);
\draw[arr] (duty) -- (rc);
\draw[arr] (rc) -- (hb);
\draw[arr] (hb) -- node[above=1pt,font=\scriptsize]{$D V_{\mathrm{bat}}$} (mot);
\draw[arr] (mot) -- (wheel);
\draw[arr] (wheel) -- (enc);
\end{tikzpicture}
\caption{Path from a velocity command to the wheels (one side). There is no
speed feedback today: the encoders feed odometry only.}
\label{fig:actuator}
\end{figure}

Figure~\ref{fig:actuator} shows the chain. \texttt{bbai\_base} receives
$(v,\omega)$, splits it into side speeds \eqref{eq:sidespeed}, and maps each
to a signed duty $D\in[-0.6,0.6]$ with the deadzone offset \eqref{eq:duty}.
librobotcontrol's \texttt{rc\_motor\_set(ch, D)} then
\begin{enumerate}
\item multiplies $D$ by the channel polarity (the fork's table is
  $\{1,-1,1,-1\}$ for M1--M4, so the wiring of each motor's leads does not
  matter to the caller),
\item sets the two direction inputs \texttt{IN1}/\texttt{IN2} of that
  H-bridge channel from the sign of $D$ (for M1 one of these is driven by the
  PRU firmware of \S\ref{sec:pru}), and
\item writes $\abs{D}$ to the channel's PWM output at \SI{25}{kHz}
  (\texttt{RC\_MOTOR\_DEFAULT\_PWM\_FREQ}).
\end{enumerate}
The \SI{25}{kHz} carrier is above hearing and far above the motor's
electrical corner frequency $R/L$, so the motor current is essentially DC
with a small ripple.

\paragraph{Drive and brake.} The TB6612FNG works in \emph{sign--magnitude}
mode: during the on-part of each PWM period the bridge applies
$\pm V_{\mathrm{bat}}$ to the motor, and during the off-part it shorts both
motor terminals together (short brake, the chip's behaviour when
\texttt{PWM} is low with a direction set). The current recirculates through
the low-side switches, so in continuous conduction the average terminal
voltage is simply
\begin{equation}\label{eq:vavg}
\bar V = D\,V_{\mathrm{bat}} - 2 i R_{\mathrm{on}},
\end{equation}
where $R_{\mathrm{on}}\approx\SI{0.5}{\ohm}$ is the on-resistance of one
switch. Setting $D=0$ (all four duty outputs zero) brakes the motors rather
than letting them coast; \texttt{rc\_motor\_free\_spin} would coast them,
but the base node does not use it.

\section{DC MOTOR MODEL}
\label{sec:motormodel}

The thesis models its permanent-magnet DC motors with the standard
electrical and mechanical equations (thesis \S5.1). Written for one motor of
this robot, with armature current $i$, motor-shaft speed $\Omega_m$ and gear
ratio $N$ (so the wheel turns at $\Omega=\Omega_m/N$),
\begin{equation}\label{eq:dcmotor}
L_a\frac{\dd i}{\dd t} = \bar V - R_a i - K_e\Omega_m,\qquad
J_m\frac{\dd\Omega_m}{\dd t} = K_t i - b_m\Omega_m - \frac{\tau_\ell}{N\eta},
\end{equation}
with armature resistance $R_a$ and inductance $L_a$, back-EMF constant
$K_e$, torque constant $K_t$ ($=K_e$ in SI units), viscous friction $b_m$,
gearbox efficiency $\eta$, and load torque $\tau_\ell$ at the wheel. The
electrical time constant $L_a/R_a$ is well under a millisecond for motors
this size, so the current settles almost instantly and
\eqref{eq:dcmotor} reduces to a first-order model for the wheel speed. With
$\tau_\ell = r F$ for a tractive force $F$ at the tyre,
\begin{equation}\label{eq:motor1st}
\tau_m\frac{\dd\Omega}{\dd t} + \Omega
 = \underbrace{\frac{K_t}{N(K_tK_e+R_ab_m)}}_{K_V} \bar V
 - \underbrace{\frac{R_a}{N^2\eta(K_tK_e+R_ab_m)}}_{K_F}\, rF,
\qquad
\tau_m=\frac{R_aJ_{\mathrm{eq}}}{K_tK_e+R_ab_m},
\end{equation}
where $J_{\mathrm{eq}}$ includes the motor rotor and the reflected inertia of
the wheel and its share of the robot's mass.

\paragraph{What this says about open-loop duty control.} In steady state
the wheel speed is
\begin{equation}\label{eq:wheelss}
v_{\mathrm{wheel}} = r\Omega = rK_V\bigl(D\,V_{\mathrm{bat}}-2iR_{\mathrm{on}}\bigr) - rK_F\,rF .
\end{equation}
The duty map \eqref{eq:duty} assumes the first term alone with a fixed
$V_{\mathrm{bat}}$. Equation \eqref{eq:wheelss} shows two errors in that
assumption:
\begin{itemize}
\item \emph{Battery voltage.} A 2S LiPo falls from \SI{8.4}{V} full to about
  \SI{7.0}{V} near empty, so the same duty gives about 17\% less speed by
  the end of a run.
\item \emph{Load.} Any force at the tyre subtracts speed in proportion to
  $K_F$. Driving straight on a level floor the load is small (rolling
  resistance), but in a skid-steer turn it is large: \eqref{eq:spinforce}
  says each side must push with about $0.42\mu_\ell W$ before the robot
  rotates at all. The second term then cancels most of the first, which is
  why low-duty turns stalled (\S\ref{sec:skidturn}).
\end{itemize}
The deadzone below about 8\% duty is the same equation with Coulomb friction
in the gearbox: the motor must make a minimum torque before it moves.

\paragraph{Identifying the model.} $K_V$ and $\tau_m$ can be measured with
the wheels lifted (where $F\approx0$) by stepping the duty and logging
\topic{/wheel\_ticks}: the final speed divided by $D V_{\mathrm{bat}}$ gives
$K_V$ (the lifted test of 2026-10-01 gives about
$\SI{0.4}{m/s}/(0.3\times\SI{7.9}{V}) \approx \SI{0.17}{m/s/V}$ at the
wheel), and the time to 63\% of it gives $\tau_m$. Repeating the step on the
floor in a spin gives the load term. Because of the coarse encoder
(\S\ref{sec:speedest}), the step response should be fitted to the
accumulated count, $\int v\,\dd t$, rather than to a differentiated speed.

\section{LOOPS ON THE ROBOT TODAY}
\label{sec:loops}

Table~\ref{tab:loops} lists every periodic loop and tracking loop in the
system, from the lowest level up.

\begin{table}[htbp]
\centering
\caption{Periodic, feedback and tracking loops}
\label{tab:loops}
\small
\begin{tabularx}{\linewidth}{>{\raggedright\arraybackslash}p{30mm}
  >{\raggedright\arraybackslash}p{19mm} >{\raggedright\arraybackslash}p{24mm} X}
\toprule
Loop & Rate & Where & What it closes \\
\midrule
PWM carrier & \SI{25}{kHz} & cape PWM & None; sets average voltage \eqref{eq:vavg}. \\
Encoder counting & per edge & eQEP $\times3$, PRU & None; hardware counters. \\
Base loop & \SI{30}{Hz} & \texttt{bbai\_base} & Reads encoders and IMU, integrates odometry \eqref{eq:odomint}, publishes \topic{/odom} and \topic{/imu/*}. \\
Command watchdog & \SI{30}{Hz} check, \SI{0.5}{s} timeout & \texttt{bbai\_base} & Stops the motors when \topic{/cmd\_vel} goes quiet. \\
Wheel speed & none & & \emph{Open loop}: duty from \eqref{eq:duty}. Proposed in \S\ref{sec:speedloop}. \\
Gyro bias (ZUPT) & \SI{30}{Hz} while still & \texttt{bbai\_base} & Tracks the gyro bias toward the measured rate when the robot is stationary (\S\ref{sec:zupt}). \\
Heading & \SI{30}{Hz} & \texttt{bbai\_base} & Complementary gyro--compass filter (\S\ref{sec:compfilt}). \\
Teleop & \SI{20}{Hz} idle, up to \SI{100}{Hz} & \texttt{bbai\_teleop} & Human in the loop. \\
Chase & \SI{15}{Hz} & \texttt{bbai\_chase} & Bearing and range to a detected dog or person (Chapter~\ref{ch:chase}). \\
Chase yaw rate & \SI{15}{Hz} & \texttt{bbai\_chase} & Gyro PI loop on the requested turn rate (\S\ref{sec:chaseyaw}; PR~\#3, untested). \\
Detection & about \SI{10}{Hz} & OAK-D-Lite & Runs on the camera; feeds the chase loop. \\
Global EKF & \SI{20}{Hz} & \texttt{ekf\_gps} & Fuses odometry, IMU yaw and GPS (Chapter~\ref{ch:fusion}). \\
SLAM & per scan (\SIrange{5}{7}{Hz}) & gmapping / hector & Scan matching against the map (Chapter~\ref{ch:slam}). \\
\bottomrule
\end{tabularx}
\end{table}

The only loops that close around the robot's motion today are the human
(teleop) and the chase controller, and both close around the motors through
the open-loop duty map. Everything between the velocity command and the
wheel is feed-forward.

\section{ENCODER SPEED ESTIMATION}
\label{sec:speedest}

Any speed loop needs a speed measurement, and the encoders are coarse: 15
counts per wheel turn, $\delta=2\pi r/15=\SI{16.6}{mm}$ of travel per count.
A side's speed averaged over its two encoders and differenced over one
\SI{30}{Hz} sample is quantized in steps of
\begin{equation}
\Delta v = \frac{\delta/2}{\Delta t} = \SI{0.25}{m/s},
\end{equation}
which is most of the robot's useful speed range. Two standard remedies, both
from the thesis's discussion of encoder sample rate (thesis \S7.5):
\begin{enumerate}
\item \emph{Fixed-window differencing.} Difference over the last $M$
  samples: $\hat v_k = (s_k-s_{k-M})/(M\Delta t)$, where $s$ is the
  accumulated distance. The quantization step falls to $\Delta v/M$ and the
  estimate lags by $M\Delta t/2$. $M=6$ (\SI{0.2}{s}) gives \SI{0.04}{m/s}
  steps and a \SI{0.1}{s} lag.
\item \emph{Period measurement.} Time-stamp each count edge and use
  $\hat v=\delta/T_{\mathrm{edge}}$. This is much better at low speed but
  needs edge timestamps, which the eQEP units can capture (their unit timer
  and capture registers) but librobotcontrol does not expose.
\end{enumerate}
Option 1 needs no driver change and is what the proposed loop below
assumes; its lag is part of the loop delay.

\section{PROPOSED: WHEEL-SPEED LOOP}
\label{sec:speedloop}

\emph{Proposed, not implemented.} A per-side PI controller around the
first-order model \eqref{eq:motor1st}, with the current duty map as
feed-forward:
\begin{equation}\label{eq:speedpi}
\begin{aligned}
e_k &= v^\ast_k - \hat v_k,\\
D_k &= \underbrace{D_{\mathrm{ff}}(v^\ast_k)}_{\eqref{eq:duty}}
      + K_p e_k + I_k,\\
I_{k+1} &= I_k + K_i\Delta t\, e_k + K_{aw}\bigl(\sat(D_k)-D_k\bigr),
\end{aligned}
\end{equation}
where $\sat$ is the $\pm0.6$ duty limit and the last term is back-calculation
anti-windup: while the output is saturated, the integrator is bled toward
the value that just reaches the limit, so it does not wind up during a
stalled turn and overshoot when the robot breaks free. The feed-forward does
most of the work on a level floor; the integrator supplies the extra duty a
skid-steer turn needs \eqref{eq:spinforce} and compensates battery sag.

\paragraph{Gain selection.} Linearising \eqref{eq:motor1st} about an
operating point, duty to side speed is $G(s)=K_D/(\tau_m s+1)$ with
$K_D=rK_VV_{\mathrm{bat}}$ (about \SI{1.3}{m/s} per unit duty from the lifted
test). With PI $C(s)=K_p+K_i/s$, choosing the PI zero on the motor pole,
$K_i/K_p=1/\tau_m$, cancels it and leaves the closed loop
\begin{equation}
\frac{\hat v}{v^\ast} = \frac{1}{\tau_c s+1},\qquad \tau_c=\frac{\tau_m}{K_pK_D}.
\end{equation}
The estimator lag $M\Delta t/2=\SI{0.1}{s}$ and the \SI{33}{ms} sample period
act as a delay $\tau_d\approx\SI{0.12}{s}$; for a phase margin of about
\SI{60}{\degree} the crossover $1/\tau_c$ should be at most about
$\pi/(6\tau_d)\approx\SI{4.4}{rad/s}$, so $\tau_c\approx\SI{0.25}{s}$. With
an assumed $\tau_m=\SI{0.1}{s}$ (to be measured as in
\S\ref{sec:motormodel}) that gives $K_p\approx0.3$ and $K_i\approx3$ per
\si{m/s}, and $K_{aw}\approx K_i/K_p$.

\paragraph{Interaction with odometry.} The loop changes nothing about how
odometry is computed. It makes the commanded and actual speeds agree, which
also makes the EKF's process model, which assumes the robot does what it is
told between measurements, more accurate.

\section{PROPOSED: YAW-RATE LOOP}
\label{sec:yawloop}

\emph{Proposed, not implemented.} Even with each side tracking its speed,
a skid-steer robot does not turn at $\omega=(v_R-v_L)/b$: the wheels slip
sideways and the effective track $\chi$ of \eqref{eq:skidkin} is larger than
$b$ and depends on the floor. The gyro measures the true yaw rate directly,
so an outer loop on yaw rate removes that uncertainty:
\begin{equation}\label{eq:yawloop}
\Delta v_k = \frac{\chi_0}{2}\,\omega^\ast_k
            + K_\omega\bigl(\omega^\ast_k-\omega_{g,k}\bigr) + J_k,\qquad
J_{k+1}=J_k+K_{\omega i}\Delta t\bigl(\omega^\ast_k-\omega_{g,k}\bigr),
\end{equation}
with $v^\ast_{L,R}=v\mp\Delta v$ passed to the speed loops, $\omega_g$ the
bias-corrected gyro rate from the base node, and $\chi_0$ a nominal effective
track (start with $b$). The gyro is sampled at the same \SI{30}{Hz} but needs
no differencing, so this loop has less delay than the speed loop and can run
somewhat faster than it. The integrator $J$ learns the floor-dependent
$\chi$; it should be frozen whenever $\abs{\omega^\ast}$ is small, so it does
not fight the straight-line drift that the heading estimate already handles.

The chase node now implements a simpler form of this loop for its own turn
commands, as a PI correction on $\omega$ ahead of the base node
(\S\ref{sec:chaseyaw}, PR~\#3, untested). Moving it into the base would
give teleop the same benefit.

\section{IMPLEMENTATION NOTES}

\begin{itemize}
\item Both loops belong in \texttt{bbai\_base}, at the existing \SI{30}{Hz},
  so they share the encoder and gyro reads already there. Raising the base
  rate to \SI{50}{Hz} would help the yaw loop more than the speed loop, whose
  limit is encoder resolution.
\item The watchdog must reset both integrators when it stops the motors, and
  the feed-forward deadzone offset should stay so that small commands still
  move the robot before the integrator catches up.
\item A parameter \texttt{closed\_loop: false} would keep the present
  behaviour available for comparison.
\end{itemize}
